% !TEX program = xelatex
\documentclass[UTF8,11pt,a4paper]{ctexart}

\usepackage[margin=25mm]{geometry}
\usepackage{amsmath,amssymb,mathtools,bm}
\usepackage{booktabs,tabularx,longtable,array,multirow}
\usepackage[ruled,linesnumbered,vlined]{algorithm2e}
\usepackage{xcolor}
\usepackage{listings}
\usepackage{enumitem}
\usepackage{hyperref}
\usepackage{float}

\hypersetup{
  colorlinks=true,
  linkcolor=blue!50!black,
  urlcolor=blue!50!black
}

\setlist[itemize]{nosep,leftmargin=2em}
\setlist[enumerate]{nosep,leftmargin=2em}
\setlength{\parskip}{0.35em}
\setlength{\parindent}{2em}

\newcolumntype{Y}{>{\raggedright\arraybackslash}X}
\newcolumntype{C}[1]{>{\centering\arraybackslash}p{#1}}
\DeclareMathOperator{\softmax}{softmax}
\DeclareMathOperator{\softplus}{softplus}
\DeclareMathOperator{\RFFT}{RFFT}
\DeclareMathOperator{\IRFFT}{IRFFT}
\DeclareMathOperator{\sg}{sg}

\lstset{
  basicstyle=\ttfamily\small,
  frame=single,
  breaklines=true,
  columns=fullflexible,
  backgroundcolor=\color{gray!5}
}

\title{MRA 异常检测模型说明手册\\
\large 基于 \texttt{mra.py} 的代码一致性科研式说明}
\author{面向仓库 \texttt{/home/akira/codespace/mra-detection} 的源码分析文档}
\date{2026年4月9日}

\begin{document}
\maketitle

\begin{abstract}
本文对仓库中的 \texttt{mra.py} 及其直接依赖模块
\texttt{freq\_reconstruction.py}、
\texttt{graph\_gcn\_reconstruction.py}、
\texttt{utils/methods/windowing.py} 等进行逐函数分析，
生成一份与代码实现严格对齐的中文 LaTeX 说明手册。
文档重点回答四类问题：
（1）频域插补模块的代码入口、训练入口与调用链；
（2）联合训练流程、损失组成、梯度传播路径与参数更新过程；
（3）由训练集得到的模型参数和统计量在测试阶段的具体作用；
（4）从数据读入到异常分数输出的完整张量流、数学原理、设计思路与实现细节。

本手册基于当前仓库默认数据和默认超参数进行验证：训练集和测试集均为
$4100\times 18$ 的多速率序列，缺失率自动推断为三组采样步长
$[1,3,5]$，对应的 \texttt{rate\_id} 为
$[0,0,0,0,0,0,1,1,1,1,1,1,2,2,2,2,2,2]$。
文档不仅描述“设计意图”，也明确指出按当前代码真实发生的行为，
例如检测器损失虽然对目标序列使用了 \texttt{detach()}，但仍会沿
\texttt{x\_complete} 这条输入路径向插补器回传梯度。
\end{abstract}

\tableofcontents
\newpage

\section{问题定义与代码入口}

\subsection{研究对象}

\texttt{mra.py} 实现的是一个“插补 + 异常检测”联合框架，面向的输入是
带结构性缺失的多变量多速率时间序列。
记原始序列为
\[
\bm{X}\in\mathbb{R}^{T\times N},
\qquad
\bm{M}^{\mathrm{str}}\in\{0,1\}^{T\times N},
\]
其中 $M^{\mathrm{str}}_{t,n}=1$ 表示第 $n$ 个变量在时刻 $t$ 缺失。
模型的目标包括两部分：
\begin{enumerate}
  \item 对缺失位置进行插补，得到完整序列 $\hat{\bm{X}}$；
  \item 对测试窗口输出异常分数并给出二值异常预测。
\end{enumerate}

该实现不是单纯的补全网络，也不是单纯的检测网络，而是
“双专家插补器 + 采样感知 Transformer 检测器”的串联结构：
\[
\text{结构缺失序列}
\rightarrow
\text{预填充}
\rightarrow
\text{GCN 专家 / 频域专家}
\rightarrow
\text{门控融合}
\rightarrow
\text{Transformer 重构检测}
\rightarrow
\text{多信号异常评分}.
\]

\subsection{主入口、训练入口与频域模块入口}

\begin{table}[H]
\centering
\caption{与本手册最相关的源码入口}
\begin{tabularx}{\textwidth}{@{}l C{1.8cm} C{1.4cm} Y@{}}
\toprule
函数或类方法 & 文件 & 起始行 & 作用 \\
\midrule
\texttt{main()} &
\texttt{mra.py} &
831 &
主程序入口；负责读数据、构窗、推断采样率、训练模型、生成补全结果和异常分数。 \\
\texttt{train\_model(...)} &
\texttt{mra.py} &
548 &
联合训练入口；频域模块、GCN 模块、门控网络和检测器都在这里一起优化。 \\
\texttt{FusionAnomalyModel.forward(...)} &
\texttt{mra.py} &
492 &
统一前向传播入口；先做插补，再送入检测器。 \\
\texttt{TwoExpertGatedImputer.forward(...)} &
\texttt{mra.py} &
298 &
双专家插补器入口；内部先调 GCN 专家，再调频域专家，再做门控融合。 \\
\texttt{FrequencyOnlyReconstructor.forward(...)} &
\texttt{freq\_reconstruction.py} &
139 &
频域专家的直接入口；先做 FFT 频谱增强，再做时域输出头修正。 \\
\texttt{FrequencyImputer.forward(...)} &
\texttt{freq\_reconstruction.py} &
111 &
频域增强核心实现；完成 \texttt{permute}、\texttt{rfft}、频率注意力、频谱增量和 \texttt{irfft}。 \\
\texttt{train\_model(...)} &
\texttt{freq\_reconstruction.py} &
153 &
独立频域重建实验的训练入口；它存在，但\textbf{不被} \texttt{mra.py} 调用。 \\
\bottomrule
\end{tabularx}
\end{table}

\subsection{从主程序到频域插补模块的真实调用链}

\begin{lstlisting}
main()
  -> train_model()
    -> model(...)
      -> FusionAnomalyModel.forward()
        -> TwoExpertGatedImputer.forward()
          -> PreImpute.fill()
          -> GraphGCNReconstructor.forward()
          -> FrequencyOnlyReconstructor.forward()
            -> FrequencyImputer.forward()
        -> SamplingAwareTransformerAD.forward()
    -> compute_losses()
    -> total_loss.backward()
    -> clip_grad_norm_(..., max_norm=1.0)
    -> AdamW.step()
\end{lstlisting}

因此，对“频域插补模块的训练入口”必须区分两种语义：
\begin{itemize}
  \item 在 \texttt{mra.py} 主方法中，频域模块\textbf{不是单独训练}，
  而是作为 \texttt{FusionAnomalyModel} 的子模块，
  通过 \texttt{mra.py::train\_model()} 与其他模块联合训练。
  \item 在 \texttt{freq\_reconstruction.py} 中确实存在一个独立的
  \texttt{train\_model()}，但那是频域重建的单模块实验脚本，
  并不参与 \texttt{mra.py} 的主流程。
\end{itemize}

\section{数据、缺失结构与多速率元数据}

\subsection{默认数据规模}

根据当前仓库中的默认数据文件：
\begin{itemize}
  \item \texttt{data/train/train\_1.csv}: $(4100,18)$；
  \item \texttt{data/test/test\_C5\_1.csv}: $(4100,18)$；
  \item \texttt{data/val/val\_1.csv}: $(1000,18)$。
\end{itemize}

第 1--6 列无缺失，第 7--12 列缺失率约为 $0.6666$，
第 13--18 列缺失率约为 $0.8$。
代码中的 \texttt{infer\_rate\_metadata()} 根据缺失率推断：
\[
\texttt{stride}=[1,1,1,1,1,1,3,3,3,3,3,3,5,5,5,5,5,5].
\]
再将不同步长映射到离散分组，得到
\[
\texttt{rate\_id}=[0,0,0,0,0,0,1,1,1,1,1,1,2,2,2,2,2,2].
\]

\begin{table}[H]
\centering
\caption{默认数据的多速率结构}
\begin{tabular}{@{}cccc@{}}
\toprule
变量索引 & 缺失率 & 推断采样步长 \texttt{stride} & 分组 \texttt{rate\_id} \\
\midrule
1--6 & 0.0 & 1 & 0 \\
7--12 & 0.6666\ (\text{约 }2/3) & 3 & 1 \\
13--18 & 0.8\ (\text{约 }4/5) & 5 & 2 \\
\bottomrule
\end{tabular}
\end{table}

\subsection{数据读取与缺失掩码}

\texttt{load\_csv\_glob\_with\_mask()} 将 CSV 读为实数矩阵并同步构造缺失掩码：
\[
\bm{X}_{\mathrm{raw}}\in\mathbb{R}^{T\times N},
\qquad
\bm{M}^{\mathrm{str}}=\mathbf{1}[\mathrm{NaN}(\bm{X}_{\mathrm{raw}})].
\]
所有文件都按行拼接，列数必须一致，否则抛错。

\subsection{仅基于观测值的标准化}

\texttt{ObservedStandardScaler} 的核心思想是：
均值和标准差只由“真实观测值”估计，而不是由填充值估计。
对第 $n$ 个变量，记观测时刻集合为
\[
\mathcal{O}_n=\{t\mid M^{\mathrm{str}}_{t,n}=0\},
\]
则
\[
\mu_n=\frac{1}{|\mathcal{O}_n|}\sum_{t\in\mathcal{O}_n}X_{t,n},
\qquad
\sigma_n=\max\!\left(
\sqrt{\frac{1}{|\mathcal{O}_n|}\sum_{t\in\mathcal{O}_n}(X_{t,n}-\mu_n)^2},
10^{-6}
\right).
\]
标准化时先把缺失值替换为列均值，再做
\[
\tilde{X}_{t,n}=\frac{X^{\mathrm{filled}}_{t,n}-\mu_n}{\sigma_n}.
\]
因此，缺失位置在标准化后恰好变成 0，这与后续
\texttt{masked\_fill(..., 0.0)} 的输入约定是一致的。

\subsection{窗口构造方式与实际数量}

\texttt{mra.py} 同时构造三类窗口：
\begin{enumerate}
  \item \textbf{训练插补窗口}
  \texttt{build\_windows()}：
  从序列起点开始逐步构造，若历史不足则用首行做前端补齐；
  默认得到 $(4100,50,18)$。
  \item \textbf{训练评分窗口}
  \texttt{build\_standard\_windows()}：
  只使用有完整 $50$ 步历史的窗口，
  默认起始索引为 99，得到 $(4001,50,18)$。
  \item \textbf{测试评分窗口}
  \texttt{build\_prompt\_test\_windows()}：
  默认只取 4000 个窗口，并用固定分界点 2000 生成标签，
  得到窗口 $(4000,50,18)$、标签 $(4000,)$。
\end{enumerate}

\begin{table}[H]
\centering
\caption{默认窗口规模}
\begin{tabular}{@{}lcc@{}}
\toprule
窗口类型 & 张量形状 & 备注 \\
\midrule
训练插补窗口 & $(4100,50,18)$ & 与原始时间步一一对应，前端允许补齐 \\
测试插补窗口 & $(4100,50,18)$ & 用于输出 \texttt{test\_completed.csv} \\
训练评分窗口 & $(4001,50,18)$ & 不做前端补齐 \\
测试评分窗口 & $(4000,50,18)$ & 对应 2000 正常 + 2000 异常标签 \\
\bottomrule
\end{tabular}
\end{table}

\section{双专家插补器的设计与实现}

\subsection{预填充 \texttt{PreImpute.fill}}

\texttt{PreImpute} 是插补器前的暖启动步骤，其输入为
\[
\bm{X}_{\mathrm{input}}\in\mathbb{R}^{B\times L\times N},
\qquad
\bm{M}^{\mathrm{model}}\in\{0,1\}^{B\times L\times N}.
\]
这里的 \texttt{model\_missing\_mask} 可能同时包含结构性缺失和训练期的随机遮挡。

预填充分两次一维扫描完成：
\begin{itemize}
  \item 正向扫描：对每个变量沿时间轴维护最近一次可见值；
  \item 反向扫描：在倒序时间上重复同样操作。
\end{itemize}

设 $x_{t,n}^{\rightarrow}$ 和 $x_{t,n}^{\leftarrow}$ 分别表示前向与后向可获得的最近观测值，
则最终种子值为
\[
x^{\mathrm{seed}}_{t,n}=
\begin{cases}
x_{t,n}, & M^{\mathrm{model}}_{t,n}=0,\\
\frac{x^{\rightarrow}_{t,n}+x^{\leftarrow}_{t,n}}{2}, &
\text{前后都可用},\\
x^{\rightarrow}_{t,n}, & \text{仅前向可用},\\
x^{\leftarrow}_{t,n}, & \text{仅后向可用},\\
0, & \text{前后都不可用}.
\end{cases}
\]

这样做的动机很直接：频域模块无法直接处理稀疏掩码，GCN 虽然能显式看掩码，
但如果输入全是 0，会放大“缺失即 0”的伪信号，因此先构造一个时域上平滑、
数值上合理的稠密种子序列。

\subsection{频域专家：\texttt{FrequencyOnlyReconstructor}}

\subsubsection{代码入口与模块组成}

在主流程中，频域专家的入口是：
\[
\texttt{TwoExpertGatedImputer.forward}
\rightarrow
\texttt{self.freq(x\_seed)}
\rightarrow
\texttt{FrequencyOnlyReconstructor.forward}
\rightarrow
\texttt{FrequencyImputer.forward}.
\]

\texttt{FrequencyOnlyReconstructor} 由两层组成：
\begin{enumerate}
  \item \texttt{FrequencyImputer}：在频谱空间做选择性增强；
  \item \texttt{output\_head}：在时域上做跨变量修正。
\end{enumerate}

\subsubsection{数学原理}

设输入为预填充后的窗口
\[
\bm{X}_{\mathrm{seed}}\in\mathbb{R}^{B\times L\times N}.
\]
代码先把时间维挪到最后一维：
\[
\bm{X}^{\top}\in\mathbb{R}^{B\times N\times L}.
\]
随后对每个变量的长度为 $L$ 的时间序列执行实数 FFT：
\[
\bm{F}=\RFFT(\bm{X}^{\top},\mathrm{dim}=2)\in\mathbb{C}^{B\times N\times K},
\qquad
K=\left\lfloor\frac{L}{2}\right\rfloor+1.
\]
默认 $L=50$，因此 $K=26$。

将复频谱拆成实部和虚部后拼接：
\[
\bm{f}=[\Re(\bm{F})\ \|\ \Im(\bm{F})]\in\mathbb{R}^{B\times N\times 2K}.
\]

然后有两条并行支路：
\begin{align}
\bm{a} &= \sigma\!\left(\mathrm{MLP}_{\mathrm{att}}(\bm{f})\right)
\in (0,1)^{B\times N\times K}, \\
[\Delta \bm{r}\ \|\ \Delta \bm{i}] &=
\mathrm{MLP}_{\mathrm{enh}}(\bm{f})
\in \mathbb{R}^{B\times N\times 2K}.
\end{align}
其中 $\bm{a}$ 是每个频率点的自适应权重，
$\Delta \bm{r}$、$\Delta \bm{i}$ 分别是对实部和虚部的可学习增量。

增强后的频谱写成
\[
\tilde{\bm{F}}
=
\bm{F}
+
\left(\bm{a}\odot \Delta \bm{r}\right)
+
j\left(\bm{a}\odot \Delta \bm{i}\right).
\]
最后执行逆变换：
\[
\tilde{\bm{X}}^{\top}=\IRFFT(\tilde{\bm{F}},n=L,\mathrm{dim}=2)
\in\mathbb{R}^{B\times N\times L},
\]
再转回
\[
\tilde{\bm{X}}\in\mathbb{R}^{B\times L\times N}.
\]

\subsubsection{为什么频域分支能工作}

频域分支的建模假设是：
多速率工业信号虽然存在缺失，但在\textbf{时间轴上仍然保留周期性、
谐波结构和低频趋势}。
相比时域逐点预测，频谱增强有三个优点：
\begin{enumerate}
  \item 周期成分在频域更稀疏，易于被一个小 MLP 选择性放大或抑制；
  \item FFT 的感受野天然覆盖整个窗口，适合捕捉长程时序模式；
  \item 当预填充已经提供粗糙但连续的种子时，
  频域网络更像是在做“全局频谱校正”而不是从零生成。
\end{enumerate}

但当前实现也有一个必须明确的约束：
\texttt{FrequencyImputer} \textbf{不直接读取掩码}，
它只看 \texttt{x\_seed}。
也就是说，缺失信息对频域分支的影响主要来自预填充结果，
以及后续门控网络对 \texttt{model\_missing\_mask} 的显式感知。

\subsubsection{输出头的作用}

\texttt{FrequencyOnlyReconstructor.forward} 在频域输出后又做了一层时域残差修正：
\[
\bm{X}_{\mathrm{freq}}
=
\tilde{\bm{X}} + \mathrm{MLP}_{\mathrm{out}}(\tilde{\bm{X}}).
\]
这里的 \texttt{Linear(num\_features, num\_features)} 是沿最后一维作用，
因此它在每个时间步做\textbf{跨变量混合}。
这一步弥补了频谱增强阶段“主要对单变量时间轨迹建模”的不足，
把变量间耦合信息补回来。

\subsubsection{张量形状追踪}

\begin{table}[H]
\centering
\caption{频域专家内部张量形状}
\begin{tabular}{@{}lll@{}}
\toprule
变量名 & 一般形状 & 默认 $L{=}50,N{=}18$ 时的末维 \\
\midrule
\texttt{x\_seed} & $(B,L,N)$ & $(B,50,18)$ \\
\texttt{x\_perm} & $(B,N,L)$ & $(B,18,50)$ \\
\texttt{xf} & $(B,N,K)$（复数） & $(B,18,26)$ \\
\texttt{feat} & $(B,N,2K)$ & $(B,18,52)$ \\
\texttt{att\_weights} & $(B,N,K)$ & $(B,18,26)$ \\
\texttt{feat\_enhanced} & $(B,N,2K)$ & $(B,18,52)$ \\
\texttt{x\_rec} & $(B,L,N)$ & $(B,50,18)$ \\
\texttt{x\_freq} & $(B,L,N)$ & $(B,50,18)$ \\
\bottomrule
\end{tabular}
\end{table}

\subsection{图专家：\texttt{GraphGCNReconstructor}}

\subsubsection{图学习器 \texttt{GraphLearner}}

图专家的输入是 \texttt{x\_seed} 和 \texttt{model\_missing\_mask}。
与频域专家不同，图专家显式使用掩码。
对每个变量 $n$，\texttt{GraphLearner} 从窗口内提取四类动态统计量：
\[
\bm{c}^{\mathrm{dyn}}_n=
\big[
\mathrm{mean}_n,\,
\mathrm{std}_n,\,
\mathrm{last}_n,\,
\mathrm{missing\_ratio}_n
\big]\in\mathbb{R}^{4}.
\]
再拼接可学习静态上下文
\[
\bm{c}^{\mathrm{static}}_n\in\mathbb{R}^{8},
\]
输入节点编码 MLP，得到
\[
\bm{h}_n\in\mathbb{R}^{d_g}.
\]

对每一对节点 $(i,j)$，代码构造四类配对特征
\[
[\bm{h}_i,\ \bm{h}_j,\ |\bm{h}_i-\bm{h}_j|,\ \bm{h}_i\odot \bm{h}_j],
\]
送入边 MLP 得到 \texttt{edge\_logits}。
此外还引入两种辅助图结构信息：
\begin{itemize}
  \item 基于可学习坐标的距离先验
  $p_{ij}=1/(1+\|\bm{p}_i-\bm{p}_j\|_2)$；
  \item 归一化节点表征之间的余弦相似度。
\end{itemize}

因此，最终邻接 logit 为
\[
\ell_{ij}
=
e_{ij}
+
\softplus(\alpha_p)\,p_{ij}
+
\softplus(\alpha_s)\,\cos(\bm{h}_i,\bm{h}_j).
\]
对角线单独加上可学习自环 logit，
再执行带温度的 \texttt{softmax}、对称化和行归一化，得到
\[
\bm{A}\in\mathbb{R}^{B\times N\times N}.
\]

\subsubsection{GCN 重构}

学到邻接矩阵后，\texttt{GraphGCNReconstructor} 做三层图卷积：
\begin{align}
\bm{H}^{(0)} &= \bm{X}_{\mathrm{seed}}\in\mathbb{R}^{B\times L\times N},\\
\bm{H}^{(1)} &= \mathrm{Dropout}\!\left(
\mathrm{LayerNorm}\!\left(
\mathrm{ReLU}(\mathrm{GCN}_{\mathrm{in}}(\bm{H}^{(0)},\bm{A}))
\right)\right),\\
\bm{H}^{(2)} &= \mathrm{Dropout}\!\left(
\mathrm{LayerNorm}\!\left(
\mathrm{ReLU}(\mathrm{GCN}_{\mathrm{mid}}(\bm{H}^{(1)},\bm{A}))
\right)\right),\\
\bm{X}_{\mathrm{gcn}} &= \mathrm{GCN}_{\mathrm{out}}(\bm{H}^{(2)},\bm{A}).
\end{align}

其中单层 GCN 的消息传递在代码中写作
\[
\texttt{einsum("bnm,bsmd->bsnd", adj, x)},
\]
其含义是：
对每个时间步 $s$，按邻接矩阵在“节点维”上做加权聚合，
而不是在时间维上做卷积。

\subsection{门控融合：\texttt{TwoExpertGatedImputer}}

\subsubsection{门控特征}

门控网络把以下 6 组标量特征沿最后一维堆叠：
\[
\big[
x_{\mathrm{seed}},
M^{\mathrm{model}},
x_{\mathrm{gcn}},
x_{\mathrm{freq}},
x_{\mathrm{gcn}}-x_{\mathrm{seed}},
x_{\mathrm{freq}}-x_{\mathrm{seed}}
\big],
\]
得到
\[
\bm{G}_{\mathrm{base}}\in\mathbb{R}^{B\times L\times N\times 6}.
\]
然后追加采样率嵌入
\[
\bm{e}_{r_n}\in\mathbb{R}^{d_r},
\]
广播后与之拼接，形成
\[
\bm{G}\in\mathbb{R}^{B\times L\times N\times (6+d_r)}.
\]
默认 $d_r=8$，因此门控输入最后一维为 14。

\subsubsection{融合公式}

门控 MLP 输出两个 logit，对应 GCN 与频域两个专家：
\[
\bm{\alpha}_{t,n}
=
\softmax(\mathrm{MLP}_{\mathrm{gate}}(\bm{g}_{t,n}))
\in\mathbb{R}^{2}.
\]
融合结果为
\[
x^{\mathrm{imputed}}_{t,n}
=
\alpha^{(\mathrm{gcn})}_{t,n}x^{\mathrm{gcn}}_{t,n}
+
\alpha^{(\mathrm{freq})}_{t,n}x^{\mathrm{freq}}_{t,n}.
\]
仅在缺失位置采用融合值：
\[
x^{\mathrm{complete}}_{t,n}
=
\begin{cases}
x^{\mathrm{input}}_{t,n}, & M^{\mathrm{model}}_{t,n}=0,\\
x^{\mathrm{imputed}}_{t,n}, & M^{\mathrm{model}}_{t,n}=1.
\end{cases}
\]

因此，门控器承担两层职责：
\begin{enumerate}
  \item 在缺失位置决定“空间相关性更可信”还是“频域周期性更可信”；
  \item 在训练时根据 \texttt{holdout} 造成的人为遮挡学习这种决策。
\end{enumerate}

\subsection{插补器输出字典}

\begin{table}[H]
\centering
\caption{双专家插补器 \texttt{outputs} 的主要键值}
\begin{tabular}{@{}lll@{}}
\toprule
键名 & 形状 & 含义 \\
\midrule
\texttt{x\_seed} & $(B,L,N)$ & 预填充结果 \\
\texttt{x\_gcn} & $(B,L,N)$ & GCN 专家输出 \\
\texttt{x\_freq} & $(B,L,N)$ & 频域专家输出 \\
\texttt{x\_imputed} & $(B,L,N)$ & 两专家加权融合值 \\
\texttt{x\_complete} & $(B,L,N)$ & 观测值保留、缺失值补全后的序列 \\
\texttt{gate\_weights} & $(B,L,N,2)$ & 两专家门控权重 \\
\texttt{adjacency} & $(B,N,N)$ & 当前批次自适应邻接矩阵 \\
\bottomrule
\end{tabular}
\end{table}

\section{采样感知 Transformer 检测器}

\subsection{输入定义}

\texttt{SamplingAwareTransformerAD} 的输入不是原始缺失序列，而是
\texttt{x\_complete}。此外还需要两类元信息：
\begin{itemize}
  \item \texttt{observed\_mask} $= 1-\texttt{structural\_missing\_mask}$；
  \item 每个变量对应的 \texttt{rate\_id} 和 \texttt{stride}。
\end{itemize}

\subsection{每变量表示构造}

对每个时间步 $t$、变量 $n$，代码拼接五类信息：
\[
\bm{v}_{t,n}
=
\big[
x^{\mathrm{complete}}_{t,n},
o_{t,n},
\bm{e}^{\mathrm{var}}_n,
\bm{e}^{\mathrm{rate}}_{r_n},
\bm{e}^{\mathrm{phase}}_{t,n}
\big].
\]
其中
\begin{itemize}
  \item $x^{\mathrm{complete}}_{t,n}$ 是补全值；
  \item $o_{t,n}$ 是结构观测标记；
  \item $\bm{e}^{\mathrm{var}}_n\in\mathbb{R}^{8}$ 是变量嵌入；
  \item $\bm{e}^{\mathrm{rate}}_{r_n}\in\mathbb{R}^{8}$ 是采样率嵌入；
  \item 相位特征定义为
  \[
  \phi_{t,n}=\frac{t\bmod s_n}{s_n}\in[0,1),
  \qquad
  \bm{e}^{\mathrm{phase}}_{t,n}
  =
  \mathrm{MLP}_{\mathrm{phase}}(\phi_{t,n}).
  \]
\end{itemize}

默认 \texttt{phase\_dim}=8，因此每变量输入维度为
\[
1+1+8+8+8=26.
\]

\subsection{Token 化与编码}

每变量输入经过 \texttt{per\_variable\_proj} 映射到
\[
d_v=\max(d/8,8)
\]
维。默认 \texttt{detector\_d\_model}=128，因此 $d_v=16$。

对每个时间步，把所有变量的三类信息拼成单个 token：
\[
\bm{z}_t=
\big[
\bm{x}^{\mathrm{complete}}_{t,:},
\bm{o}_{t,:},
\mathrm{reshape}(\bm{u}_{t,:})
\big]
\in\mathbb{R}^{2N+Nd_v}.
\]
在默认 $N=18,d_v=16$ 下，
\[
2N+Nd_v=2\times 18 + 18\times 16 = 324.
\]
再用 \texttt{token\_proj} 投影到 $d=128$ 维，
叠加标准正弦位置编码后进入 Transformer Encoder。

\subsection{重构输出}

编码器输出经 \texttt{reconstruction\_head} 生成增量
\[
\bm{\Delta}\in\mathbb{R}^{B\times L\times N},
\]
最终输出为
\[
\bm{X}_{\mathrm{det}}
=
\bm{X}_{\mathrm{complete}}+\bm{\Delta}.
\]
代码还保存逐元素绝对误差
\[
\bm{E}=|\bm{X}_{\mathrm{det}}-\bm{X}_{\mathrm{complete}}|.
\]

这意味着检测器不是从零生成一个新序列，
而是在补全序列上学习一个“残差偏移”。
正常训练数据上，较理想的状态是 $\bm{\Delta}$ 小；
异常窗口上，$\bm{\Delta}$ 往往变大。

\section{联合训练流程与参数更新机制}

\subsection{训练期的三类掩码}

训练循环每次从 DataLoader 取出
\[
\bm{X}_{\mathrm{true}},\ \bm{M}^{\mathrm{str}}
\in\mathbb{R}^{B\times L\times N}.
\]
随后通过 \texttt{sample\_holdout\_mask()}，
在\textbf{原本可观测的位置}上随机制造一层额外遮挡：
\[
\bm{M}^{\mathrm{hold}}\in\{0,1\}^{B\times L\times N}.
\]
模型实际看到的掩码为
\[
\bm{M}^{\mathrm{model}}
=
\max(\bm{M}^{\mathrm{str}},\bm{M}^{\mathrm{hold}}).
\]
输入张量是
\[
\bm{X}_{\mathrm{input}}
=
\bm{X}_{\mathrm{true}}\odot(1-\bm{M}^{\mathrm{model}}).
\]

这个设计的意义是：
\begin{itemize}
  \item 原始结构缺失保持不变；
  \item 额外随机遮挡为插补器制造监督目标；
  \item 至少保证每个 batch 样本若存在观测值，则会有至少一个点被 holdout。
\end{itemize}

\subsection{损失函数}

公共的掩码 MSE 定义为
\[
\mathcal{L}_{\mathrm{mask}}(\hat{\bm{X}},\bm{X},\bm{M})
=
\frac{\sum_{b,t,n}\big((\hat{X}_{b,t,n}-X_{b,t,n})M_{b,t,n}\big)^2}
{\max\left(\sum_{b,t,n}M_{b,t,n},1\right)}.
\]

当前代码中的五项损失是
\begin{align}
\mathcal{L}_{\mathrm{fusion}}
&=
\mathcal{L}_{\mathrm{mask}}
(\bm{X}_{\mathrm{imputed}},\bm{X}_{\mathrm{true}},\bm{M}^{\mathrm{hold}}),\\
\mathcal{L}_{\mathrm{gcn}}
&=
\mathcal{L}_{\mathrm{mask}}
(\bm{X}_{\mathrm{gcn}},\bm{X}_{\mathrm{true}},\bm{M}^{\mathrm{hold}}),\\
\mathcal{L}_{\mathrm{freq}}
&=
\mathcal{L}_{\mathrm{mask}}
(\bm{X}_{\mathrm{freq}},\bm{X}_{\mathrm{true}},\bm{M}^{\mathrm{hold}}),\\
\mathcal{L}_{\mathrm{graph}}
&=
\frac{1}{BN}\sum_{b,n}(A_{b,n,n}-\tau_{\mathrm{diag}})^2,\\
\mathcal{L}_{\mathrm{det}}
&=
\mathcal{L}_{\mathrm{mask}}
\big(
\bm{X}_{\mathrm{det}},
\sg(\bm{X}_{\mathrm{complete}}),
1-\bm{M}^{\mathrm{str}}
\big).
\end{align}

总损失为
\[
\mathcal{L}
=
1.0\,\mathcal{L}_{\mathrm{fusion}}
+
0.4\,\mathcal{L}_{\mathrm{gcn}}
+
0.4\,\mathcal{L}_{\mathrm{freq}}
+
0.5\,\mathcal{L}_{\mathrm{det}}
+
0.1\,\mathcal{L}_{\mathrm{graph}}.
\]

\subsection{参数分组}

为了说明“谁被什么损失更新”，把参数粗分为四组：
\begin{align}
\theta_G &: \text{GraphLearner 与三层 GCN 的参数},\\
\theta_F &: \text{频域模块 \texttt{attention}、\texttt{freq\_enhance}、\texttt{output\_head}},\\
\theta_M &: \text{门控 MLP 与插补器中的 \texttt{rate\_embedding}},\\
\theta_D &: \text{检测器全部参数}.
\end{align}

\subsection{梯度路由：代码真实发生的更新关系}

\begin{table}[H]
\centering
\caption{各损失项实际更新哪些参数}
\begin{tabularx}{\textwidth}{@{}l Y Y@{}}
\toprule
损失项 & 直接依赖的中间量 & 会被更新的参数 \\
\midrule
$\mathcal{L}_{\mathrm{fusion}}$ &
\texttt{x\_imputed} &
$\theta_G,\theta_F,\theta_M$ \\
$\mathcal{L}_{\mathrm{gcn}}$ &
\texttt{x\_gcn} &
$\theta_G$ \\
$\mathcal{L}_{\mathrm{freq}}$ &
\texttt{x\_freq} &
$\theta_F$ \\
$\mathcal{L}_{\mathrm{graph}}$ &
\texttt{adjacency} &
仅 GraphLearner 中生成邻接矩阵的参数，属于 $\theta_G$ 的子集 \\
$\mathcal{L}_{\mathrm{det}}$ &
\texttt{detector(outputs["x\_complete"])} &
$\theta_D$，以及所有会影响 \texttt{x\_complete} 的插补器参数，即 $\theta_G,\theta_F,\theta_M$ \\
\bottomrule
\end{tabularx}
\end{table}

\paragraph{关键说明 1：\texttt{detach()} 并没有把检测器和插补器完全隔离。}
当前代码中
\[
\texttt{detector\_target = outputs["x\_complete"].detach()}
\]
只对\textbf{目标分支}停止梯度，但检测器前向仍然是
\[
\texttt{self.detector(outputs["x\_complete"], ...)}.
\]
因此 $\mathcal{L}_{\mathrm{det}}$ 依然会沿
\texttt{detector\_reconstruction -> detector -> x\_complete}
这条路径回传到插补器。
在本仓库默认配置下，对 \texttt{detector\_loss} 单独做反向传播时，
GCN、频域、门控、采样率嵌入和检测器参数都会得到非零梯度。

\paragraph{关键说明 2：检测器训练目标包含 holdout 位置的“伪观测”。}
\texttt{detector\_loss} 使用的掩码是
\[
1-\bm{M}^{\mathrm{str}},
\]
而不是
\[
1-\bm{M}^{\mathrm{model}}.
\]
所以在训练时，人为 holdout 的位置仍被当作“观测位置”参与检测器损失；
但这些位置的目标值已经不是原始真值，而是
\texttt{x\_complete.detach()} 中的插补结果。
换言之，检测器在这部分位置上学习的是“重构插补器输出”，而不是“重构原始观测值”。

\subsection{参数更新公式}

优化器为 AdamW，代码顺序为：
\begin{enumerate}
  \item \texttt{optimizer.zero\_grad()}
  \item \texttt{loss.backward()}
  \item \texttt{clip\_grad\_norm\_(..., max\_norm=1.0)}
  \item \texttt{optimizer.step()}
\end{enumerate}

若将裁剪后的梯度记为 $\hat{\bm{g}}_t$，则 AdamW 的参数更新可写作
\begin{align}
\bm{m}_t &= \beta_1\bm{m}_{t-1} + (1-\beta_1)\hat{\bm{g}}_t,\\
\bm{v}_t &= \beta_2\bm{v}_{t-1} + (1-\beta_2)\hat{\bm{g}}_t^{\odot 2},\\
\theta_t &=
\theta_{t-1}
-\eta
\frac{\bm{m}_t/(1-\beta_1^t)}
{\sqrt{\bm{v}_t/(1-\beta_2^t)}+\epsilon}
-\eta\lambda \theta_{t-1}.
\end{align}
其中默认超参数是
\[
\eta=10^{-3},\qquad \lambda=10^{-4}.
\]

\subsection{训练伪代码}

\begin{lstlisting}
for epoch in range(E):
    for x_true, structural_missing_mask in loader:
        holdout_mask = sample_holdout_mask(structural_missing_mask, holdout_ratio)
        model_missing_mask = max(structural_missing_mask, holdout_mask)
        x_input = x_true * (1 - model_missing_mask)

        outputs = model(
            x_input,
            model_missing_mask,
            structural_missing_mask,
            rate_id,
            stride,
        )
        losses = compute_losses(outputs, ...)
        total_loss = weighted_sum(losses)

        optimizer.zero_grad()
        total_loss.backward()
        clip_grad_norm_(model.parameters(), 1.0)
        optimizer.step()
\end{lstlisting}

\section{训练集派生量在测试阶段的作用}

\subsection{需要区分两类“参数”}

用户提出“训练集更新的参数和在测试集中的作用”，在当前代码里至少包含两种不同来源：
\begin{enumerate}
  \item \textbf{通过梯度下降学到的模型参数}：
  例如 GCN 权重、频域 MLP、门控网络、Transformer 权重；
  \item \textbf{由训练集统计估计得到的非梯度量}：
  例如标准化均值方差、\texttt{rate\_id}/\texttt{stride}、阈值、分数归一化统计量。
\end{enumerate}

\subsection{梯度学习得到的参数}

\begin{table}[H]
\centering
\caption{模型参数在测试阶段的具体作用}
\begin{tabularx}{\textwidth}{@{}l Y@{}}
\toprule
参数组 & 在测试阶段的作用 \\
\midrule
GraphLearner 参数 &
根据测试窗口的统计特征动态生成邻接矩阵，决定节点间如何传播信息。 \\
GCN 参数 &
沿邻接图把其他变量的信息聚合到当前变量，用于补全缺失值。 \\
频域模块参数 &
对测试窗口执行 FFT 频谱增强，恢复周期、趋势和谐波结构。 \\
门控网络参数 &
根据当前位置的缺失状态、两个专家的输出及采样率嵌入，自适应分配 GCN/频域权重。 \\
插补器中的 \texttt{rate\_embedding} &
把采样率分组编码成向量，影响门控决策。 \\
检测器中的变量嵌入、采样率嵌入、相位 MLP &
让 Transformer 区分“不同变量”“不同采样率”和“周期内不同相位”的正常模式。 \\
Transformer Encoder 与重构头 &
对补全后的测试序列做重构；偏离训练正常模式的窗口会产生更大的重构误差。 \\
\bottomrule
\end{tabularx}
\end{table}

\subsection{由训练集估计的非梯度量}

\begin{table}[H]
\centering
\caption{训练集派生统计量及其测试作用}
\begin{tabularx}{\textwidth}{@{}l Y Y@{}}
\toprule
统计量 & 训练阶段如何得到 & 测试阶段如何使用 \\
\midrule
\texttt{scaler.mean\_}, \texttt{scaler.scale\_} &
对训练集真实观测值逐列估计均值和标准差 &
用于把测试集映射到与训练集相同的标准化坐标系；输出补全结果时再逆变换。 \\
\texttt{rate\_id}, \texttt{stride} &
根据训练集缺失率推断不同采样步长与采样率组 &
作为门控网络和检测器的元信息输入；当前实现默认直接沿用于测试集。 \\
训练分数均值/标准差 &
\texttt{normalize\_scores()} 对训练分数求均值和标准差 &
用于把测试分数做 Z-score 标准化，保证 train/test 在同一评分尺度。 \\
训练分布偏移特征均值/标准差 &
\texttt{compute\_distribution\_shift\_scores()} 对训练补全窗口统计特征做估计 &
衡量测试补全窗口离训练分布有多远。 \\
阈值 $\theta$ &
由平滑后的训练分数通过“高斯阈值与分位数阈值取大”得到 &
用于把测试异常分数转换为二值预测。 \\
\bottomrule
\end{tabularx}
\end{table}

\paragraph{一个实现上的现实问题。}
\texttt{model.pt} 保存了模型参数、\texttt{rate\_id}、\texttt{stride} 和损失权重，
但\textbf{没有保存} \texttt{ObservedStandardScaler} 的均值方差，
也没有保存用于 \texttt{normalize\_scores()} 和
\texttt{compute\_distribution\_shift\_scores()} 的全部训练统计量。
因此，如果未来要把 \texttt{model.pt} 单独拿出来做部署，
还需要额外保存或重建这些训练集统计量。

\section{测试推理、序列补全与异常评分}

\subsection{序列补全}

\texttt{reconstruct\_full\_sequence()} 在测试阶段使用
\[
\bm{M}^{\mathrm{model}}=\bm{M}^{\mathrm{str}},
\]
即不再注入 holdout。
函数对所有插补窗口逐批前向传播，但只取每个窗口的最后一个时间步：
\[
\hat{\bm{x}}_i
=
\bm{X}^{\mathrm{complete}}_{i,L-1,:}\in\mathbb{R}^{N}.
\]
由于训练和测试插补窗口都是 $(4100,50,18)$，
最终拼接后会得到长度正好为 4100 的补全序列，
并保存为 \texttt{train\_completed.csv} 与 \texttt{test\_completed.csv}。

\subsection{窗口统计量收集}

\texttt{collect\_window\_statistics()} 对评分窗口收集三种统计量：
\begin{enumerate}
  \item 检测器重构误差分数
  \[
  s^{\mathrm{det}}_i
  =
  \frac{\sum_{t,n}
  \big((X^{\mathrm{det}}_{i,t,n}-X^{\mathrm{complete}}_{i,t,n})\,O_{i,t,n}\big)^2}
  {\max(\sum_{t,n}O_{i,t,n},1)};
  \]
  \item 专家分歧度
  \[
  s^{\mathrm{dis}}_i
  =
  \frac{1}{LN}\sum_{t,n}|x^{\mathrm{gcn}}_{i,t,n}-x^{\mathrm{freq}}_{i,t,n}|;
  \]
  \item 门控熵
  \[
  s^{\mathrm{ent}}_i
  =
  -\frac{1}{LN}\sum_{t,n}\sum_{k=1}^{2}\alpha^{(k)}_{i,t,n}\log\alpha^{(k)}_{i,t,n}.
  \]
\end{enumerate}

其中门控熵会被保存，但\textbf{不参与最终评分融合}。

\subsection{分布偏移分数}

\texttt{compute\_distribution\_shift\_scores()} 不是直接在原始窗口上工作，
而是在“补全后的训练窗口”和“补全后的测试窗口”上提取统计特征。
它调用 \texttt{build\_window\_feature\_matrix()}，后者固定使用四个列块：
\[
[0{:}6],\ [6{:}12],\ [12{:}18],\ [0{:}18].
\]
对每个列块分别提取均值和标准差，再额外追加
\begin{itemize}
  \item 一阶差分绝对值均值；
  \item 前 5 个非直流 FFT 幅值均值。
\end{itemize}

在当前 $N=18$ 的实现下，特征维度可精确写为
\[
D=2\times (6+6+6+18)+18+18=108.
\]
随后对这些 108 维统计特征做训练分布标准化，
并输出 RMS 型偏移分数：
\[
s^{\mathrm{shift}}_i=
\sqrt{\frac{1}{D}\sum_{d=1}^{D}
\left(\frac{f_{i,d}-\mu^{\mathrm{train}}_d}{\sigma^{\mathrm{train}}_d}\right)^2 }.
\]

\subsection{分数归一化与融合}

三类分数都先用训练集自身的均值方差标准化：
\[
\tilde{s}_i=
\frac{s_i-\mu_{\mathrm{train}}}{\max(\sigma_{\mathrm{train}},10^{-6})}.
\]
最终原始异常分数为
\[
a^{\mathrm{raw}}_i=
|\tilde{s}^{\mathrm{det}}_i|
+
0.25\,\tilde{s}^{\mathrm{dis}}_i
+
1.0\,\tilde{s}^{\mathrm{shift}}_i.
\]

\subsection{平滑、阈值与最终标签}

\texttt{apply\_ewaf\_by\_segments()} 用 EWAF 做平滑：
\[
a^{\mathrm{smooth}}_1=a^{\mathrm{raw}}_1,\qquad
a^{\mathrm{smooth}}_i=
\alpha a^{\mathrm{raw}}_i+(1-\alpha)a^{\mathrm{smooth}}_{i-1},
\]
默认 $\alpha=0.3$。

阈值由
\[
\theta=
\max\left(
\mu_{\mathrm{train}} + 2.0\,\sigma_{\mathrm{train}},
Q_{0.95}(\bm{a}^{\mathrm{smooth}}_{\mathrm{train}})
\right)
\]
给出，再产生原始预测
\[
\hat{y}^{\mathrm{raw}}_i=\mathbf{1}[a^{\mathrm{smooth}}_i\ge \theta].
\]
最后 \texttt{apply\_min\_anomaly\_duration()} 要求异常片段长度至少为 50，
短于 50 的连续异常片段会被抑制为 0。

\section{完整张量流与默认形状总表}

\begin{longtable}{@{}p{3.0cm} p{3.7cm} p{5.0cm} p{2.4cm}@{}}
\caption{从输入到输出的完整张量流}\\
\toprule
阶段 & 变量 & 说明 & 默认形状 \\
\midrule
\endfirsthead
\toprule
阶段 & 变量 & 说明 & 默认形状 \\
\midrule
\endhead
\bottomrule
\endfoot
数据读取 &
$\bm{X}_{\mathrm{raw}}$ &
原始 CSV 拼接后的实值矩阵 &
$(4100,18)$ \\
数据读取 &
$\bm{M}^{\mathrm{str}}$ &
结构缺失掩码 &
$(4100,18)$ \\
标准化 &
$\bm{X}_{\mathrm{scaled}}$ &
以训练观测统计量标准化后的矩阵 &
$(4100,18)$ \\
训练插补构窗 &
\texttt{train\_impute\_windows} &
\texttt{build\_windows()} 的输出 &
$(4100,50,18)$ \\
训练评分构窗 &
\texttt{train\_eval\_windows} &
\texttt{build\_standard\_windows()} 的输出 &
$(4001,50,18)$ \\
测试评分构窗 &
\texttt{test\_eval\_windows} &
\texttt{build\_prompt\_test\_windows()} 的输出 &
$(4000,50,18)$ \\
批输入 &
\texttt{x\_true} &
训练 batch 的真实窗口 &
$(B,50,18)$ \\
批输入 &
\texttt{structural\_missing\_mask} &
训练 batch 的结构掩码 &
$(B,50,18)$ \\
训练增强 &
\texttt{holdout\_mask} &
随机从观测点采样的附加遮挡 &
$(B,50,18)$ \\
模型输入 &
\texttt{x\_input} &
被 \texttt{model\_missing\_mask} 遮掉后的输入 &
$(B,50,18)$ \\
预填充 &
\texttt{x\_seed} &
双向最近观测传播的暖启动结果 &
$(B,50,18)$ \\
图学习 &
\texttt{adjacency} &
批相关自适应邻接矩阵 &
$(B,18,18)$ \\
GCN 专家 &
\texttt{x\_gcn} &
图卷积专家重构值 &
$(B,50,18)$ \\
频域预处理 &
\texttt{x\_perm} &
\texttt{permute(0,2,1)} 后的频域输入 &
$(B,18,50)$ \\
频域 FFT &
\texttt{xf} &
实数 FFT 频谱，复数张量 &
$(B,18,26)$ \\
频域特征 &
\texttt{feat} &
实部虚部拼接后的实数特征 &
$(B,18,52)$ \\
频域专家 &
\texttt{x\_freq} &
频域增强 + 输出头修正后的结果 &
$(B,50,18)$ \\
门控输入 &
\texttt{gate\_input} &
六类标量特征 + 采样率嵌入 &
$(B,50,18,14)$ \\
门控输出 &
\texttt{gate\_weights} &
两个专家的 softmax 权重 &
$(B,50,18,2)$ \\
插补输出 &
\texttt{x\_imputed} &
两专家融合值 &
$(B,50,18)$ \\
补全输出 &
\texttt{x\_complete} &
保留观测值后的完整序列 &
$(B,50,18)$ \\
检测器输入 &
\texttt{per\_variable\_input} &
值、观测标记、变量嵌入、采样率嵌入、相位嵌入 &
$(B,50,18,26)$ \\
每变量投影 &
\texttt{rate\_aware\_features} &
\texttt{per\_variable\_proj} 的输出 &
$(B,50,18,16)$ \\
Token 化 &
\texttt{token\_input} &
按时间步拼成单 token &
$(B,50,324)$ \\
Transformer 编码 &
\texttt{encoded} &
Encoder 输出 &
$(B,50,128)$ \\
检测器输出 &
\texttt{detector\_reconstruction} &
残差重构后的检测器输出 &
$(B,50,18)$ \\
误差张量 &
\texttt{detector\_error} &
逐元素绝对误差 &
$(B,50,18)$ \\
窗口统计 &
\texttt{detector\_score} &
每个窗口 1 个重构误差分数 &
$(B,)$ \\
窗口统计 &
\texttt{disagreement\_score} &
每个窗口 1 个专家分歧度 &
$(B,)$ \\
窗口统计 &
\texttt{gate\_entropy} &
每个窗口 1 个门控熵 &
$(B,)$ \\
序列重构输出 &
\texttt{train\_complete\_scaled} &
只取每个窗口最后一步并拼接 &
$(4100,18)$ \\
测试预测 &
\texttt{test\_scores} &
EWAF 平滑后的测试异常分数 &
$(4000,)$ \\
测试预测 &
\texttt{final\_prediction} &
二值异常预测 &
$(4000,)$ \\
\end{longtable}

\section{代码层面的实现约束、假设与注意事项}

\subsection{分布偏移评分硬编码了 18 维特征布局}

\texttt{build\_window\_feature\_matrix()} 固定切片
\texttt{slice(0,6)}、\texttt{slice(6,12)}、\texttt{slice(12,18)}、
\texttt{slice(0,18)}。
这意味着当前偏移评分逻辑默认变量数至少能与这种三组结构对齐；
如果换数据集，尤其是变量数量或分组方式变化，
这一部分需要同步改写，否则特征工程语义会失真。

\subsection{测试标签来自固定分割，而不是文件中真实标签}

\texttt{build\_prompt\_test\_windows()} 并不读取外部标签文件，
而是把前 2000 个窗口记为正常、后 2000 个窗口记为异常。
这说明当前评估脚本默认测试集已经按“前正常、后异常”的时间顺序组织。
如果测试集组织方式改变，评价指标会直接失效。

\subsection{\texttt{fdr} 字段在代码里等于召回率}

\texttt{compute\_binary\_classification\_metrics()} 中
\texttt{fdr} 的实现是
\[
\texttt{tp}/(\texttt{tp}+\texttt{fn}),
\]
它与 \texttt{recall} 完全相同。
从命名上看这并不符合通常统计学中对 FDR 的定义，
因此解读 \texttt{metrics.json} 时要以源码实现为准。

\subsection{默认参数量}

在默认超参数
\[
N=18,\ L=50,\ d=128,\ d_g=32,\ d_{\mathrm{gcn}}=32,\ d_r=8
\]
下，\texttt{FusionAnomalyModel} 总参数量约为 693260。
这是一个中等规模模型，适合在当前数据规模上做窗口级训练，
但如果未来窗口更长、变量更多，Transformer 和门控分支的开销会明显上涨。

\section{结论}

从源码角度看，\texttt{mra.py} 的方法学可以概括为：
\begin{enumerate}
  \item 先利用 \texttt{PreImpute} 把缺失窗口变成稠密种子；
  \item 再用 GCN 专家建模变量间结构相关性，用频域专家建模时间周期结构；
  \item 用门控网络在位置级别自适应融合两类重构；
  \item 用采样感知 Transformer 对补全序列做残差式重构检测；
  \item 最后结合检测误差、专家分歧和分布偏移三种信号输出异常分数。
\end{enumerate}

频域插补模块的\textbf{代码入口}是
\texttt{mra.py::TwoExpertGatedImputer.forward()} 中的
\texttt{self.freq(x\_seed)}，
其\textbf{训练入口}是 \texttt{mra.py::train\_model()}，
而不是单独的优化器或单独的训练脚本。

从训练机制上看，当前实现最值得注意的地方有两点：
\begin{enumerate}
  \item 检测器损失会反向影响插补器，因而这是\textbf{真正的耦合式联合训练}；
  \item 检测器在训练期会把 holdout 位置上的插补值当作重构目标的一部分。
\end{enumerate}

如果后续要把该方法整理成正式论文或进一步工程化部署，
建议优先补齐三类工作：
\begin{enumerate}
  \item 显式保存标准化器和训练评分统计量；
  \item 把分布偏移特征工程从硬编码列切片改为可配置分组；
  \item 明确决定检测器是否应当继续对插补器回传梯度，
  若不希望耦合，应在检测器输入侧对 \texttt{x\_complete} 再做一次 \texttt{detach()}。
\end{enumerate}

\end{document}
